% GMC: Graph Multi-view Completion for Drug Repositioning
% Bioinformatics (OUP) submission
%
% Template: oup-authoring-template v1.3 (2026/05/28)
%   - modern:  Bioinformatics journal house style (required)
%   - webpdf:  hyperlink-aware PDF output
%   - large:   standard page dimensions

% ===== flushend prevention (MUST come before \documentclass) =====
\makeatletter
\expandafter\def\csname ver@flushend.sty\endcsname{9999/99/99}
\makeatother

\documentclass[webpdf,modern,large]{oup-authoring-template}

% --- OUP template metadata ---------------------------------------------------
\journaltitle{[Journal Name]}
\pubyear{[Year]}
\copyrightyear{[Year]}
\vol{[Vol]}
\issue{[Issue]}
\DOI{[DOI to be added]}

\usepackage{graphicx}
\usepackage{booktabs}
\usepackage{amsmath}
\usepackage{amssymb}
\newcounter{algline}
\usepackage{multirow}
\newcommand{\gmcmode}[1]{\emph{mode~#1}}   % ablation remove-one-at-a-time mode label
\usepackage{xcolor}
\usepackage{hyperref}

% --- Template bug workarounds (cls v1.3) ---
\makeatletter
\providecommand{\authormark}[1]{}
\long\def\abstract#1{\let\subheading\abstsubheading\global\def\@abstract{#1}}
\let\flushend\relax
\def\societylogo{}
\makeatother


\newcommand{\R}{\mathbb{R}}
\newcommand{\E}{\mathbb{E}}
\newcommand{\tr}{\mathrm{tr}}

\begin{document}

\title{GMC: Graph Multi-view Completion for Drug Repositioning}

\author{[Author Name(s)]}
\address{[Affiliation, Department, University, City, Country]}
\corresp{[Corresponding author: email@example.com]}

\date{}

\abstract{%
\textbf{Motivation:} Drug repositioning finds new therapeutic indications for approved drugs at lower cost and risk than de novo discovery. Under the random-entry masking protocol (CVa) used by the multi-similarity drug--disease benchmarks, prediction is a \emph{matrix completion} problem and the leading methods all rely on global low-rank factorization, which has two gaps: it drops the local neighborhood signal that label-propagation methods use, and it merges the multi-similarity views inside a single objective so that errors in one view contaminate the whole reconstruction rather than being corrected by the others.\par
\textbf{Results:} We introduce \textbf{GMC (Graph Multi-view Completion)}, which couples a cold-start-restricted KNN initialization with scale-free multi-view rank fusion within one nuclear-norm completion framework: the same similarity inputs drive a matrix geometry on the joint bipartite+similarity block and a tensor geometry that keeps each similarity slice separate. Under 10-fold CVa, GMC attains the highest AUPR on all four benchmarks, significantly ahead of every published baseline on CTDdataset2023 and Ydataset ($p=0.002$, paired fold-level Wilcoxon signed-rank test); on Fdataset and Cdataset, where the block geometry alone already attains the full-model AUPR, it leads the closest competitor, DNMFDDA, by margins that do not reach significance ($p=0.131$). Every score is a completion score: with all known associations removed GMC collapses to the base rate on every dataset, so it propagates observed associations rather than inferring pairs from similarity alone.\par
\textbf{Availability:} Source code and datasets are available at [GitHub URL: to be added upon publication].\par
\textbf{Contact:} [Corresponding author email: to be added]%
}
\maketitle

\section{Introduction}

Bringing a new drug to market costs over \$2.6 billion and takes 10--15 years \citep{dimasi2003,paul2010}; repurposing compounds that have already cleared safety testing avoids much of that cost and risk \citep{chong2007}. Computational repurposing methods differ mainly in the structural prior they place on the drug--disease association matrix, ranging from network propagation \citep{yamanishi2008,martinez2015,luo2016} and low-rank factorization \citep{dai2015,ezzat2017,zheng2013,yang2020msbmf} to graph neural networks \citep{fu2022,gao2023} and matrix completion \citep{yang2019,luo2018,mongia2022}.

A common thread across the strongest of these methods is the use of \emph{multiple} drug and disease similarity views at once. The multi-similarity family supplies five drug similarities (chemical structure, ATC codes, side effects, drug--drug interactions, target profiles) and two disease similarities (phenotype, disease ontology) for the Fdataset, Cdataset, CTDdataset2023, and Ydataset benchmarks \citep{yang2020msbmf}, and the methods that fuse these views, MSBMF \citep{yang2020msbmf}, HGIMC \citep{yang2021hgimc}, and the multi-similarity geometric matrix factorization (multiGMF) \citep{yang2026multigmf}, consistently outperform single-view predictors.

Which paradigm wins depends on the evaluation protocol. Under disease-centric cross-validation (CVc), entire held-out diseases are unobserved and candidates are ranked within the full pool, a setting where local label-propagation and graph-filtering methods do well. Under random-entry masking (CVa), each fold holds out 10\% of the positive associations and an equal number of negatives, zeroing those entries and asking the model to recover them from a partially observed matrix; this is a \emph{matrix completion} problem \citep{candes2009}, and the methods that lead it are all global low-rank factorizations or completions: DNMFDDA (deep semi-NMF), ITRPCA (incomplete tensor robust PCA), OMC (overlap matrix completion), and BNNR (nuclear-norm completion). We adopt CVa because it is the standard setting of the multi-similarity benchmark suite, including multiGMF. Random-entry masking also matches the practical situation in which some known drug--disease associations are unobserved at prediction time, whereas CVc targets fully cold-start novel diseases and answers a different question.

The pure completion paradigm has two weaknesses that motivate our design. First, a low-rank projection is global, so it discards the \emph{local} neighborhood structure on which label propagation relies: two drugs sharing identical low-rank coordinates are treated the same even if only one sits next to a known disease in the similarity graph. Second, the 5+2 similarity views are merged inside a single objective, so a systematic error in one view spreads through the whole reconstruction instead of being corrected by the others.

We propose \textbf{GMC} (Graph Multi-view Completion), a new model built on two principles:

\begin{enumerate}
    \item \textbf{Cold-start-restricted fill (initialization of the completion).} All-zero rows and columns (novel drugs or diseases) carry no signal for a rank solver, so GMC seeds exactly those positions with the similarity-weighted average of the labels of the $k=10$ nearest neighbors over the similarity graphs \citep{yang2019omc2}. Partially observed entities keep their sparse zeros, so the completion propagates from true labels only. This is the local, graph-based warm start that pure completion lacks: an initialization of the solver, not a standalone predictor.

    \item \textbf{Multi-view low-rank completion with scale-free fusion.} A single nuclear-norm completion is solved in two geometries of the \emph{same} similarity data: a matrix geometry on the joint bipartite+similarity block $\left[\begin{smallmatrix}\mathbf{W}_{dd} & \mathbf{F}\\ \mathbf{F}^\top & \mathbf{W}_{rr}\end{smallmatrix}\right]$ (denoising mode), and a tensor geometry over the per-similarity tensors. The two readouts live on different scales, so they are rank-normalized to $[0,1]$ before combination; their errors are decorrelated, so the fusion adds robustness precisely where the two views differ most (Section~\ref{sec:results}).
\end{enumerate}

GMC is a single low-rank completion prior, instantiated in two geometries of the same similarity data and combined into one estimator, with the fill acting as the solver's initialization on novel rows and columns. The ablation study (Section~\ref{sec:results}) isolates each ingredient's contribution.

We evaluate GMC on the four benchmarks under 10-fold CVa. GMC achieves the highest mean AUPR on every dataset (Figure~\ref{fig:main}, Table~\ref{tab:main_results}): significantly ahead of all published baselines on CTDdataset2023 and Ydataset ($p=0.002$, paired fold-level Wilcoxon signed-rank test), and ahead of the closest competitor DNMFDDA on Fdataset and Cdataset by margins that do not reach significance. We state this asymmetry explicitly rather than as a footnote: on Fdataset and Cdataset the block geometry alone already attains the full-model AUPR, so no fusion headroom remains to turn into a significant gap, whereas on CTDdataset2023 and Ydataset the block geometry underfits and the tensor geometry has genuine work to do. This is a property of the data --- how much of the recoverable signal the block geometry already holds on its own --- not of the model, and the ablation (Section~\ref{sec:ablation}) localizes the asymmetry to the tensor view. We claim a demonstrated advantage only where the data supports one. A limit that bounds all of these claims: GMC's signal enters only through the observed entries of the association matrix. The similarity views define the geometry in which completion happens --- they are not an independent source of labels --- and with every known association removed the model collapses to the base rate (Section~\ref{sec:case}). The CVa results are therefore a completion result, not evidence that drug--disease associations can be inferred from similarity alone.

\section{Materials and Methods}

\subsection{Datasets and Similarities}

We evaluate on four drug--disease association benchmarks from the multi-similarity literature \citep{yang2020msbmf}. Each dataset provides an association matrix $\mathbf{A} \in \{0,1\}^{n_d \times n_s}$ (diseases $\times$ drugs), together with five drug--drug similarity matrices and two disease--disease similarity matrices. The drug similarities are chemical structure (ChemS, Tanimoto over CDK fingerprints \citep{steinbeck2003}), ATC code (AtcS), side effects (SideS), drug--drug interactions (DDIS), and target profiles (TargetS); the disease similarities are phenotype (PhS, MimMiner over MeSH descriptions \citep{vandriel2006}) and disease ontology (DoS). These 5+2 similarity views are not part of the original published association benchmarks; they were computed by \citet{yang2020msbmf} from public databases (DrugBank, SIDER, OMIM, and the disease ontology). Table~\ref{tab:datasets} summarizes the datasets. We use the mean-fused similarity views $\mathbf{W}_{rr} = \frac{1}{5}\sum_{m=1}^5 \mathbf{S}^{(m)}_{rr}$ and $\mathbf{W}_{dd} = \frac{1}{2}\sum_{m=1}^2 \mathbf{S}^{(m)}_{dd}$ for the block completion, while the tensor view uses the per-similarity matrices directly.

\begin{table*}[tb]
\centering
\caption{Dataset statistics. All four datasets provide the full 5+2 multi-similarity views.}
\label{tab:datasets}
\begin{tabular}{lcccc}
\toprule
Dataset & Diseases ($n_d$) & Drugs ($n_s$) & Known associations & Density \\
\midrule
Fdataset & 313 & 593 & 1,933 & $1.04 \times 10^{-2}$ \\
Cdataset & 409 & 663 & 2,532 & $9.34 \times 10^{-3}$ \\
CTDdataset2023 & 278 & 1,237 & 3,740 & $1.09 \times 10^{-2}$ \\
Ydataset & 655 & 1,478 & 8,448 & $8.73 \times 10^{-3}$ \\
\bottomrule
\end{tabular}
\end{table*}

All four datasets fall in the moderate-density regime ($\rho \approx 10^{-2}$).

\subsection{GMC: Method Description}
\label{sec:method}

Random-entry masking turns drug repositioning into a matrix-completion problem: recover the held-out associations from the masked training matrix $\mathbf{M}$ (held-out entries zeroed). In this regime the decisive prior is low rank: true associations lie on a low-dimensional manifold of drug--disease geometry, which is exactly the capacity the completion family (BNNR, ITRPCA, OMC) contributes. GMC is a single completion model built on that prior. It embeds the masked matrix in the drug--disease similarity space and solves one nuclear-norm minimization, fitted only on the observed entries, in two complementary geometries of the \emph{same} similarity data: a matrix geometry that fuses all similarity views into one joint block, and a tensor geometry that keeps each similarity view a separate slice. The two solutions estimate the same completion from the same evidence, so they are combined on a scale-free axis at fixed weights into one estimator. A cold-start-restricted fill initializes the solver on degenerate all-zero rows and columns without disturbing the observed entries, and training entries are restored verbatim at the output.

\subsubsection{Cold-Start-Restricted Fill (Initialization of the Completion)}

A rank-based solver needs a starting value at every entry, but all-zero rows and columns of the masked matrix $\mathbf{M} \in \R^{n_d \times n_s}$ carry no gradient signal, so the completion is degenerate there. GMC seeds exactly those positions with a similarity-weighted neighbor prior \citep{yang2019omc2}: for each all-zero row (novel disease) $i$, the fill averages the association profiles of the $k=10$ most similar diseases using the fused disease similarity $\mathbf{W}_{dd}$:

\begin{equation}
\tilde{\mathbf{M}}_{ic} = \frac{\sum_{j \in \mathcal{N}_k(i)} \mathbf{W}_{dd}(i,j)\,\mathbf{M}_{jc}}{\sum_{j' \in \mathcal{N}_k(i)} \mathbf{W}_{dd}(i,j')},
\label{eq:knnfill}
\end{equation}

where $\mathcal{N}_k(i)$ is the set of the $k{=}10$ nearest neighbors of $i$ (self excluded), taken over diseases for the row fill and over drugs for the column fill; the all-zero column fill is identical in form, with $\mathbf{W}_{dd}\!\to\!\mathbf{W}_{rr}$ and $\mathcal{N}_k(i)\!\to\!\mathcal{N}_k(c)$. The fill is applied only where the row or column carries \emph{no} observed associations:

\begin{equation}
\mathbf{F}_{ij} = \begin{cases}
\tilde{\mathbf{M}}_{ij} & \text{if row } i \text{ or column } j \text{ is all-zero in } \mathbf{M} \\
\mathbf{M}_{ij} & \text{otherwise}
\end{cases}
\label{eq:coldstart}
\end{equation}

Rows/columns with any known label keep their sparse zeros, so the observed structure reaches the solver unchanged and the completion propagates exclusively from true labels for those entities, while genuinely novel drugs/diseases still receive a data-driven prior. This restriction is what keeps the fill an initialization rather than a predictor: an unrestricted fill overwrites the known--unknown boundary and weakens the completion, and in our experiments the restricted fill consistently improves over the full fill.

\subsubsection{Matrix Geometry: Nuclear Norm over the Fused Joint Block}

The first geometry fuses all similarity views into one weighted graph and completes the association matrix in the single low-rank space this graph defines. The completed matrix is read as the top-right block of the low-rank denoising of the joint structure matrix
\begin{equation}
\mathbf{T} = \begin{bmatrix} \mathbf{W}_{dd} & \mathbf{F} \\ \mathbf{F}^\top & \mathbf{W}_{rr} \end{bmatrix} \in \R^{(n_d+n_s)\times(n_d+n_s)},
\label{eq:block}
\end{equation}
obtained by solving the bounded nuclear-norm minimization \citep{yang2019}
\begin{equation}
\begin{split}
\min_{\mathbf{X}} \; &\tfrac{1}{2}\lVert \mathcal{P}_\Omega(\mathbf{X}-\mathbf{T}) \rVert_F^2 + \alpha \lVert \mathbf{X} \rVert_*, \\
&\text{subject to } \mathbf{X}\in[0,1],
\end{split}
\label{eq:bnnr}
\end{equation}
via ADMM \citep{boyd2011} with singular-value soft-thresholding \citep{cai2010} (rank-capped at $400$, chosen as the center of the sensitivity grid (Section~\ref{sec:param}) and sufficient for all four datasets), where $\Omega$ is the set of observed (non-zero) entries of $\mathbf{T}$: only the true labels and the similarity blocks are fitted, so the low-rank projection is free to move the zero-valued (unknown) positions rather than being pinned to the fill. Because the similarity blocks are dense and observed, this is close to the \emph{denoising} mode of the block completion: the similarity structure regularizes the low-rank projection of the fill. The drug--disease block is recovered as $\mathbf{M}_{\text{block}} = \mathbf{X}_{[0:n_d,\,n_d:n_d+n_s]}$ and clipped to $[0,1]$. Treating the similarity structure as part of the low-rank block, rather than as a separate constraint, is what lets the completion borrow strength across the joint graph. A symmetric block with rank cap $400$ and the observed-nonzero mask is used (the rank cap matters on the larger datasets, whose association matrices have higher intrinsic rank). This is the matrix geometry of the rank prior; the next subsection lifts the same prior to a tensor geometry that keeps the individual similarity views instead of fusing them.

\subsubsection{Tensor Geometry: Nuclear Norm over the Per-Modality Slices}

The second geometry lifts the same rank prior to third order: instead of fusing the similarity views, it keeps each one a separate slice and solves the completion with incomplete tensor robust PCA \citep{yang2023itrpc} over the per-similarity tensors $\boldsymbol{\mathcal{T}}_{rr} \in \R^{n_s\times n_s\times 5}$ and $\boldsymbol{\mathcal{T}}_{dd} \in \R^{n_d\times n_d\times 2}$. The tensor solver proceeds by the alternating direction method of multipliers in the Fourier domain,
\begin{equation}
\begin{split}
\min_{\boldsymbol{\mathcal{L}},\boldsymbol{\mathcal{S}}} \;
&\sum_{k=1}^{5} \lambda_k^{(r)} \lVert \boldsymbol{\mathcal{L}}^{(k)}_{rr} \rVert_{\oplus,p} \\
&{}+ \sum_{k=1}^{2} \lambda_k^{(d)} \lVert \boldsymbol{\mathcal{L}}^{(k)}_{dd} \rVert_{\oplus,p} + \lVert \boldsymbol{\mathcal{S}} \rVert_1, \\
&\boldsymbol{\mathcal{P}}_\Omega\big(\boldsymbol{\mathcal{L}}+\boldsymbol{\mathcal{S}}\big) = \boldsymbol{\mathcal{P}}_\Omega\big(\mathbf{M}\big),
\end{split}
\label{eq:itrpca}
\end{equation}
where $\lVert\cdot\rVert_{\oplus,p}$ is the tensor (weighted Schatten-$p$) nuclear norm in the Fourier (t-SVD) domain and $\Omega$ is the set of observed entries; the WKNN fill ($k=30$, $r=0.95$; the solver's own internal fill, distinct from GMC's cold-start fill) and the two tensors are the method's standard inputs. The low-rank component gives the tensor readout $\mathbf{M}_{\text{ten}}$. The two geometries encode the same evidence differently: the matrix geometry uses only the \emph{fused} similarities, while the tensor geometry keeps each similarity slice separate, so its low-rank basis can represent view-specific structure that averaging destroys. The geometries are complementary rather than redundant.

\subsubsection{Combining the Two Geometries into One Estimator}

The two geometries produce two estimates of the \emph{same} masked completion from the \emph{same} similarity data, though they present the association evidence differently: the matrix geometry embeds the cold-start-filled matrix $\mathbf{F}$ in the joint block, whereas the tensor geometry operates on the bare masked matrix $\mathbf{M}$, with the solver's own internal WKNN fill. Both readouts serve a single estimator of the same masked completion. Because the two readouts live on different scales, a raw weighted sum would be dominated by whichever has larger magnitude, so each is rank-normalized to $[0,1]$ (global value rank) before combining:
\begin{equation}
\mathbf{F}_{\text{fuse}} = \frac{w_{\text{b}}\,\mathrm{rnorm01}(\mathbf{M}_{\text{block}}) + w_{\text{t}}\,\mathrm{rnorm01}(\mathbf{M}_{\text{ten}})}{w_{\text{b}}+w_{\text{t}}},
\label{eq:fusion}
\end{equation}
with equal weights $w_{\text{b}}=w_{\text{t}}=0.5$, yielding a single estimator. Because the two geometries have decorrelated errors (the matrix readout draws on the joint space, the tensor readout on per-modality structure), the combined estimator is at least as accurate as either readout alone within fold-level noise, and strictly better where the geometries differ most (Section~\ref{sec:results}). The rank-fused scores are returned directly, with training entries restored verbatim:
\begin{equation}
\hat{\mathbf{M}}_{ij} = \begin{cases}
\mathbf{M}_{ij} & \text{if } \mathbf{M}_{ij} \neq 0 \\
(\mathbf{F}_{\text{fuse}})_{ij} & \text{otherwise}
\end{cases}
\label{eq:restore}
\end{equation}

\subsubsection{Relationship to the Completion and Propagation Families}

GMC is a single model designed from one low-rank completion prior. It solves one nuclear-norm minimization in two complementary geometries of the same similarity data, and the local similarity-graph structure enters that single prior through the cold-start-restricted fill and the block geometry. Pure completion ignores local neighborhood structure (its low-rank projection is global), while pure label propagation lacks the global low-rank capacity to recover randomly masked entries (it only smooths what is already there). GMC addresses both within one estimator (the low-rank prior supplies the global recovery, and the fill and block geometry supply the local structure as part of the same completion objective), which is what lets it achieve its accuracy under CVa.

\subsubsection{Algorithm}

Algorithm~\ref{alg:gmc} summarizes the complete procedure, and Figure~\ref{fig:framework} provides an overview of the GMC model. The block-completion solver is based on the bounded nuclear-norm method \citep{yang2019} and the tensor-completion solver is based on the tensor robust PCA method \citep{yang2023itrpc}; both are re-implemented in Python, and their published MATLAB implementations are used for the baselines, so all GMC components run entirely in Python.

\begin{algorithm}[htbp]
\caption{GMC: Graph Multi-view Completion}
\label{alg:gmc}
\small
\setlength{\tabcolsep}{0pt}
\begin{tabular}{@{\stepcounter{algline}\thealgline.\quad}p{0.55\linewidth}@{\quad}p{0.20\linewidth}@{}}
\textbf{Require:} fused $\mathbf{W}_{rr}, \mathbf{W}_{dd}$; per-similarity tensors $\{\mathbf{S}^{(m)}_{rr}\}, \{\mathbf{S}^{(m)}_{dd}\}$; masked $\mathbf{M}$ & \\
\textbf{Ensure:} $\hat{\mathbf{M}}$ (completed association matrix) & \\
$\mathbf{F} \gets \mathrm{KNNcold}(\mathbf{M}, \mathbf{W}_{rr}, \mathbf{W}_{dd})$ & Eq.~\eqref{eq:knnfill},~\eqref{eq:coldstart} \\
$\mathbf{T} \gets \mathrm{block}(\mathbf{W}_{dd}, \mathbf{F}, \mathbf{W}_{rr})$ & Eq.~\eqref{eq:block} \\
$\mathbf{M}_{\text{block}} \gets \mathrm{BlockComplete}(\mathbf{T}, \Omega{=}\text{observed mask})$ & Eq.~\eqref{eq:bnnr} \\
$\mathbf{M}_{\text{ten}} \gets \mathrm{TensorComplete}(\{\mathbf{S}^{(m)}_{rr}\}, \{\mathbf{S}^{(m)}_{dd}\}, \mathbf{M})$ & Eq.~\eqref{eq:itrpca} \\
$\mathbf{F}_{\text{fuse}} \gets \tfrac{1}{2}\,\mathrm{rnorm01}(\mathbf{M}_{\text{block}}) + \tfrac{1}{2}\,\mathrm{rnorm01}(\mathbf{M}_{\text{ten}})$ & Eq.~\eqref{eq:fusion} \\
$\hat{\mathbf{M}} \gets \mathbf{M} + (\mathbf{1} - \mathbf{1}_{\mathbf{M} \neq 0}) \odot \mathbf{F}_{\text{fuse}}$ & Eq.~\eqref{eq:restore} \\
\textbf{return} $\hat{\mathbf{M}}$ & \\
\end{tabular}
\end{algorithm}

\begin{figure*}[tb]
\centering
\includegraphics[width=\textwidth]{figures/fig1_framework}
\caption{Overview of the GMC model. (a) The five drug and two disease similarity views are mean-fused into $\mathbf{W}_{rr}$ and $\mathbf{W}_{dd}$. (b) A cold-start-restricted fill seeds the completion solver: similarity-weighted neighbor averages are injected only into all-zero rows and columns of the masked matrix $\mathbf{M}$, yielding the initialization $\mathbf{F}$ without disturbing the observed entries. (c) GMC solves a single nuclear-norm completion in two geometries of the same similarity data: the matrix geometry completes the joint block $\left[\begin{smallmatrix}\mathbf{W}_{dd} & \mathbf{F}\\ \mathbf{F}^{\top} & \mathbf{W}_{rr}\end{smallmatrix}\right]$ (rank-capped ADMM, observed-nonzero mask), and the tensor geometry completes the per-similarity tensors. The two readouts of the same masked completion are rank-normalized ($\mathrm{rnorm01}$) and combined at fixed equal weights into one estimator, used directly (no post-hoc filter).}
\label{fig:framework}
\end{figure*}

\subsection{Evaluation Protocol}
\label{sec:eval}

We evaluate using 10-fold random-entry cross-validation (CVa), the protocol of the multi-similarity family: for each fold, 10\% of the positive associations (and an equal number of negatives) are held out, their entries are zeroed, and the task is to recover them. This is a matrix-completion setting. We report the area under the receiver-operating characteristic curve (AUROC), the area under the precision--recall curve (AUPR), and top-$K$ precision at $K = 10, 20$. Given extreme class imbalance (positive rate $\le 1.1\%$ in all datasets), AUPR is the primary metric \citep{saito2015}. All experiments use a fixed seed (12345); folds are identical across methods, so per-fold paired comparisons are valid.

GMC's hyperparameters (fill, block structure, rank cap, observation mask, and fusion weights) were selected on validation folds independent of the reported test folds: candidate settings were screened on a fresh independent CVa split (seed 24680), and the selected settings were then confirmed on the reported test folds. The two fold sets agree to within $\approx 0.001$--$0.003$ AUPR on every dataset (Fdataset $0.0010$, Cdataset $0.0033$, CTDdataset2023 $0.0021$, Ydataset $0.0008$), so the settings are not overfit to either fold set. Settings were frozen before the reported folds were scored, so the reported test folds did not participate in selection.

Comparison methods are run with the authors' own code under the identical folds. We compare against nine published methods: BNNR (nuclear-norm completion, ChemS+PhS) \citep{yang2019}, HGIMC (ChemS+PhS and fused 5+2) \citep{yang2021hgimc}, MSBMF \citep{yang2020msbmf}, DDA-SKF (similarity kernel fusion + LapRLS) \citep{gao2022ddaskf}, OMC (overlap matrix completion with KNN cold-start fill) \citep{yang2019omc2}, ITRPCA (incomplete tensor robust PCA) \citep{yang2023itrpc}, DNMFDDA (deep semi-NMF) \citep{yang2025dnmf}, multiGMF (both the ChemS+PhS and full 5+2 configurations) \citep{yang2026multigmf}, and NMF-DR (network-based non-negative matrix factorization, an independent method from an unrelated group) \citep{sadeghi2022}. AUPR differences between the model and each baseline are assessed with two-sided paired fold-level Wilcoxon signed-rank tests; the ten folds are paired by construction because every method is scored on the same held-out entries. (For $n=10$ folds, $p=0.002$ is the smallest two-sided $p$-value attainable, corresponding to all ten folds differing in the same direction.) To control the false-discovery rate across the eleven published baseline configurations per dataset, we additionally apply a Benjamini--Hochberg correction to the AUPR comparisons.

\section{Results}
\label{sec:results}

\subsection{Main Results}

Table~\ref{tab:main_results} and Figure~\ref{fig:main} compare GMC against the nine published methods on all four datasets. GMC attains the highest mean AUPR on all four datasets (Table~\ref{tab:main_results}): significantly ahead of every published baseline on CTDdataset2023 and Ydataset ($p=0.002$), and leading the closest competitor DNMFDDA on Fdataset and Cdataset without reaching significance ($p=0.131$ on both). After Benjamini--Hochberg correction across the eleven published baseline configurations per dataset, GMC remains significantly ahead of every baseline on CTDdataset2023 and Ydataset ($11/11$) and of all but DNMFDDA on Fdataset and Cdataset ($10/11$). The dataset dependence is not a statistical residual we set aside: it is explained and verified below (third observation) and by the ablation (Section~\ref{sec:ablation}).

\begin{table*}[tb]
\centering
\caption{Comparison with published methods (mean $\pm$ std, 10-fold CVa). AUPR is the primary metric; the best AUPR per column among the proposed method and baselines is in \textbf{bold}. OMC / ITRPCA / DNMFDDA use the fused 5+2 input where their pipelines accept it; (C+Ph) denotes the ChemS+PhS input configuration. NMF-DR \citep{sadeghi2022} is an independent, non-Yang baseline run from the authors' code.}
\label{tab:main_results}
\small
\begin{tabular}{llcccc}
\toprule
Dataset & Method & AUROC & AUPR & P@10 & P@20 \\
\midrule
\multirow{12}{*}{Fdataset}
& BNNR & $0.9321 \pm 0.0161$ & $0.6002 \pm 0.0244$ & 0.99 & 0.99 \\
& OMC & $0.9447 \pm 0.0120$ & $0.6106 \pm 0.0271$ & 0.99 & 0.99 \\
& ITRPCA & $0.9521 \pm 0.0105$ & $0.6321 \pm 0.0241$ & 1.00 & 1.00 \\
& DNMFDDA & $0.9623 \pm 0.0064$ & $0.6453 \pm 0.0313$ & 1.00 & 0.97 \\
& HGIMC (C+Ph) & $0.9232 \pm 0.0204$ & $0.5501 \pm 0.0319$ & 0.98 & 0.98 \\
& HGIMC & $0.9252 \pm 0.0170$ & $0.5517 \pm 0.0304$ & 0.98 & 0.98 \\
& MSBMF & $0.9477 \pm 0.0126$ & $0.6086 \pm 0.0345$ & 0.95 & 0.96 \\
& DDA-SKF & $0.9223 \pm 0.0145$ & $0.4055 \pm 0.0449$ & 0.87 & 0.83 \\
& NMF-DR & $0.6915 \pm 0.0185$ & $0.0550 \pm 0.0125$ & 0.50 & 0.36 \\
& multiGMF (C+Ph) & $0.9308 \pm 0.0195$ & $0.5800 \pm 0.0272$ & 1.00 & 1.00 \\
& multiGMF & $0.9534 \pm 0.0110$ & $0.6253 \pm 0.0255$ & 1.00 & 0.99 \\
& GMC (ours) & $0.9588 \pm 0.0096$ & \textbf{$0.6569 \pm 0.0253$} & 0.99 & 0.98 \\
\midrule
\multirow{12}{*}{Cdataset}
& BNNR & $0.9502 \pm 0.0117$ & $0.6805 \pm 0.0354$ & 0.99 & 0.99 \\
& OMC & $0.9587 \pm 0.0089$ & $0.6900 \pm 0.0351$ & 0.99 & 0.99 \\
& ITRPCA & $0.9661 \pm 0.0086$ & $0.7038 \pm 0.0226$ & 1.00 & 1.00 \\
& DNMFDDA & $0.9727 \pm 0.0053$ & $0.7221 \pm 0.0174$ & 1.00 & 1.00 \\
& HGIMC (C+Ph) & $0.9439 \pm 0.0110$ & $0.6369 \pm 0.0238$ & 0.99 & 0.99 \\
& HGIMC & $0.9484 \pm 0.0128$ & $0.6387 \pm 0.0231$ & 0.99 & 0.99 \\
& MSBMF & $0.9628 \pm 0.0103$ & $0.6761 \pm 0.0381$ & 0.91 & 0.94 \\
& DDA-SKF & $0.9328 \pm 0.0100$ & $0.4558 \pm 0.0302$ & 0.89 & 0.87 \\
& NMF-DR & $0.7007 \pm 0.0114$ & $0.0350 \pm 0.0074$ & 0.25 & 0.20 \\
& multiGMF (C+Ph) & $0.9490 \pm 0.0114$ & $0.6630 \pm 0.0197$ & 1.00 & 1.00 \\
& multiGMF & $0.9673 \pm 0.0072$ & $0.6931 \pm 0.0210$ & 1.00 & 1.00 \\
& GMC (ours) & $0.9703 \pm 0.0064$ & \textbf{$0.7285 \pm 0.0222$} & 1.00 & 1.00 \\
\midrule
\multirow{12}{*}{CTDdataset2023}
& BNNR & $0.8629 \pm 0.0087$ & $0.2339 \pm 0.0229$ & 0.83 & 0.72 \\
& OMC & $0.8830 \pm 0.0095$ & $0.2637 \pm 0.0310$ & 0.85 & 0.73 \\
& ITRPCA & $0.9065 \pm 0.0104$ & $0.3287 \pm 0.0248$ & 0.81 & 0.83 \\
& DNMFDDA & $0.9114 \pm 0.0083$ & $0.2906 \pm 0.0253$ & 0.90 & 0.82 \\
& HGIMC (C+Ph) & $0.8588 \pm 0.0123$ & $0.2581 \pm 0.0259$ & 0.84 & 0.78 \\
& HGIMC & $0.8761 \pm 0.0121$ & $0.2601 \pm 0.0255$ & 0.80 & 0.78 \\
& MSBMF & $0.8982 \pm 0.0065$ & $0.2817 \pm 0.0355$ & 0.78 & 0.73 \\
& DDA-SKF & $0.8959 \pm 0.0075$ & $0.2746 \pm 0.0262$ & 0.83 & 0.79 \\
& NMF-DR & $0.6986 \pm 0.0125$ & $0.0715 \pm 0.0105$ & 0.46 & 0.40 \\
& multiGMF (C+Ph) & $0.8858 \pm 0.0104$ & $0.2582 \pm 0.0209$ & 0.78 & 0.78 \\
& multiGMF & $0.9104 \pm 0.0059$ & $0.3184 \pm 0.0296$ & 0.84 & 0.81 \\
& GMC (ours) & $0.9147 \pm 0.0087$ & \textbf{$0.3714 \pm 0.0269$} & 0.90 & 0.86 \\
\midrule
\multirow{12}{*}{Ydataset}
& BNNR & $0.9515 \pm 0.0032$ & $0.7214 \pm 0.0130$ & 1.00 & 1.00 \\
& OMC & $0.9599 \pm 0.0029$ & $0.7279 \pm 0.0133$ & 1.00 & 1.00 \\
& ITRPCA & $0.9643 \pm 0.0037$ & $0.7086 \pm 0.0133$ & 1.00 & 1.00 \\
& DNMFDDA & $0.9710 \pm 0.0025$ & $0.7251 \pm 0.0137$ & 1.00 & 1.00 \\
& HGIMC (C+Ph) & $0.9535 \pm 0.0026$ & $0.7035 \pm 0.0113$ & 0.99 & 0.99 \\
& HGIMC & $0.9605 \pm 0.0046$ & $0.7048 \pm 0.0114$ & 0.99 & 0.99 \\
& MSBMF & $0.9661 \pm 0.0036$ & $0.7219 \pm 0.0164$ & 0.92 & 0.95 \\
& DDA-SKF & $0.9409 \pm 0.0045$ & $0.3466 \pm 0.0234$ & 0.85 & 0.86 \\
& NMF-DR & $0.6823 \pm 0.0061$ & $0.0331 \pm 0.0065$ & 0.45 & 0.35 \\
& multiGMF (C+Ph) & $0.9568 \pm 0.0033$ & $0.6876 \pm 0.0155$ & 1.00 & 1.00 \\
& multiGMF & $0.9669 \pm 0.0042$ & $0.7101 \pm 0.0151$ & 1.00 & 1.00 \\
& GMC (ours) & $0.9690 \pm 0.0026$ & \textbf{$0.7404 \pm 0.0124$} & 1.00 & 1.00 \\
\bottomrule
\end{tabular}
\end{table*}

Three observations stand out. First, under this protocol the winner is the completion family, not the propagation family: the completion baselines (BNNR, OMC, ITRPCA, DNMFDDA, MSBMF) dominate the label-propagation and diffusion methods (multiGMF, HGIMC, DDA-SKF) on every dataset, and NMF-DR, which diffuses associations through the similarity graphs and factorizes the diffused matrix, attains the lowest AUPR (0.03--0.07) while its AUROC remains $\approx 0.69$ (Table~\ref{tab:main_results}): its transductive diffusion scores are dominated by neighborhood association density, so under the $\approx 1{:}95$ imbalance its held-out positives reach only the mid-ranks, and AUPR, set by the extreme top of the ranking, collapses. This confirms that random-entry masking is a low-rank completion problem, not a propagation problem. Second, GMC beats the entire completion family: it couples global low-rank capacity with the cold-start graph prior, and the tensor$+$block rank fusion completes the model where the block view alone is weakest (Section~\ref{sec:ablation}): on Fdataset and Cdataset the block view alone already attains the full-model AUPR, while the tensor view is decisive on CTDdataset2023 ($+0.045$ AUPR, $p=0.002$) and adds $+0.011$ on Ydataset.

Third, the size of GMC's advantage tracks how much fusion headroom the data leaves. On Fdataset and Cdataset the tensor view is neutral, so GMC reduces in effect to its block view --- and the closest competitor DNMFDDA receives the same fused 5+2 similarities that the block view regularizes, so the two work from the same substrate and differ mainly in the factorization geometry applied to it. GMC's mean advantage there is $0.012$ AUPR on Fdataset and $0.006$ on Cdataset: smaller than the fold-level standard deviation of either model, with fold-level differences changing sign (GMC is ahead in 6 of the 10 folds on each), so ten folds cannot separate it from noise. On CTDdataset2023 the block view is instead the weakest of GMC's two readouts, and on Ydataset it is still the dominant component but leaves headroom for the tensor view as well (Figure~\ref{fig:ablation}). In both cases the tensor view has real work to do and the gap becomes significant. We do not claim a demonstrated advantage on Fdataset and Cdataset; what is demonstrated there is that GMC is no worse than the best baseline, and the ablation shows the asymmetry is carried by the tensor view rather than by any component failing on those datasets.

\begin{figure*}[tb]
\centering
\includegraphics[width=\textwidth]{figures/fig2_main}
\caption{AUPR of GMC vs.\ published baselines under 10-fold CVa (mean $\pm$ std over the ten folds; one panel per dataset). GMC (proposed method) attains the highest mean AUPR on all four datasets.}
\label{fig:main}
\end{figure*}

\subsection{Ablation Study}
\label{sec:ablation}

Figure~\ref{fig:ablation} isolates each ingredient of GMC on all four datasets using a \emph{remove-one-module} design on the fresh validation folds (the folds on which the configuration was selected, so every mode is scored on the identical held-out entries; the fresh-fold \gmcmode{0} differs from the reported test-fold numbers (Table~\ref{tab:main_results}) by $\approx 0.001$--$0.003$ AUPR, and all gaps below are measured on these fresh folds). Three components make up the model: a cold-start KNN fill, a block completion view, and a tensor completion view. Each \gmcmode{X} starts from the full reported configuration (\gmcmode{0}) and ablates exactly one component: \gmcmode{A} drops the cold-start KNN fill ($\mathrm{fill}=\mathrm{none}$), \gmcmode{B} drops the block completion view ($w_{\text{b}}=0$, tensor-only), and \gmcmode{C} drops the tensor view ($w_{\text{t}}=0$, block-only). The drop from \gmcmode{0} to a given mode quantifies that component's contribution. The scale-free rank fusion is a numerical detail rather than a fourth component and is examined separately below.

% Ablation data are shown in Figure~\ref{fig:ablation}; key per-component gaps are reported in the narrative below.

\begin{figure*}[tb]
\centering
\includegraphics[width=0.72\textwidth]{figures/fig3_ablation}
\caption{Remove-one-module ablation of GMC (AUPR, fresh validation folds). Four panels, one per dataset, each showing the full model (\gmcmode{0}) and the three modes with a single architectural component removed: \gmcmode{A} the cold-start fill, \gmcmode{B} the block completion view, \gmcmode{C} the tensor view. The block view is the engine on Fdataset, Cdataset, and Ydataset (the \gmcmode{B} bar is the lowest there); the tensor view is decisive on CTDdataset2023 (and moderate on Ydataset), where the block view alone is weakest; the fill is load-bearing on every dataset.}
\label{fig:ablation}
\end{figure*}

The block completion carries the model wherever low-rank recovery dominates. Removing it (\gmcmode{B}) is the largest single drop on Fdataset ($0.073$), Cdataset ($0.070$), and Ydataset ($0.120$), all $p=0.002$; CTDdataset2023 is the exception, where its $0.039$ drop is smaller than the tensor's ($0.045$) and the fill's ($0.042$), the single dataset in which the block is not the dominant component. On Fdataset and Cdataset the block-with-fill variant alone reaches the full-model AUPR ($0.6572$ vs.\ $0.6579$; $0.7259$ vs.\ $0.7252$). The tensor view pays its way only where the block geometry underfits: ablating it (\gmcmode{C}) is neutral on Fdataset and Cdataset ($-0.001$/$+0.001$, $p>0.6$) but costs $0.045$ AUPR on CTDdataset2023 and $0.011$ on Ydataset (both $p=0.002$). The cold-start fill is load-bearing on every dataset (\gmcmode{A}, $p\le0.014$; per-dataset gaps in Figure~\ref{fig:ablation}), most on CTDdataset2023, where it ties the tensor view as the largest per-mode loss: all-zero rows and columns carry no label evidence, so without the neighbor-derived prior the low-rank completion has nothing to constrain or pull from at exactly those positions.

The scale-free rank fusion is a numerical detail rather than a fourth component: it lets the block and tensor readouts, which sit on different scales, be combined at fixed equal weights. We tested this choice directly by substituting the raw weighted sum for the rank-normalized one; AUPR changes by $|\Delta|\le0.0003$ on every dataset (all $p>0.15$). At the operating point the combined score is invariant to how the two readouts are normalized before the fixed-weight combination.

Two further controls from the unified-config grid refine this reading. Both are evaluated at fixed tensor-heavy weights ($w_{\text{b}}=0.2$, $w_{\text{t}}=0.8$) as a diagnostic, which makes their reported effects conservative: the observed-mask and rank-cap choices act on the block view, and at $0.2$ block weight any block-side effect on the fused score is attenuated relative to the reported equal weights. Constraining only the observed (non-zero) block entries rather than the full block pays off where the tensor view leads: it gains $0.007$ AUPR on CTDdataset2023 and $0.021$ on Ydataset at these fixed weights. The rank cap is not sensitive either: moving it from 200 to 400 changes AUPR by $0.008$ on CTDdataset2023 and $-0.006$ on Ydataset at these weights, and the parameter analysis (Figure~\ref{fig:param}) re-examines the choice at the reported equal weights; the $400$ budget is a conservative default rather than a tuned setting.

\subsection{Parameter Sensitivity}
\label{sec:param}

GMC uses the hyperparameter set $\alpha=0.5$, 40 iterations, rank cap 400, $k=10$ cold-start KNN fill, and equal ($0.5/0.5$) block/tensor fusion weights. To check whether the reported numbers depend on fragile tuning, Figure~\ref{fig:param} perturbs each hyperparameter one at a time around this center, holding the rest fixed, on an \emph{equidistant 5-point grid} over the fresh validation folds (the folds on which the settings were selected): $\alpha \in \{0.1,0.3,0.5,0.7,0.9\}$, maxiter $\in \{20,30,40,50,60\}$, rank cap $\in \{200,300,400,500,600\}$, $k \in \{5,10,15,20,25\}$, and $w_t \in \{0.1,0.3,0.5,0.7,0.9\}$ (with $w_{\text{b}}=1-w_t$), all five axes on all four datasets. Across the interior of the grid GMC is flat: every setting that does not suppress one of the components deviates from the center by at most $0.018$ AUPR (CTDdataset2023 $0.3514$ at 20 iterations; the next largest is $0.016$ at $w_t=0.3$), within the fold-level standard-deviation band ($0.010$--$0.028$), so the reported results do not sit on a sharp peak. The informative exceptions are the settings that eliminate a component. A nearly absent nuclear-norm term ($\alpha=0.1$) costs $0.044$ AUPR on Fdataset and $0.036$ on Cdataset ($p=0.002$), the two datasets where the block view does most of the work. The one-sided fusion weights, in turn, ablate a completion view and mirror the remove-one-module result (Figure~\ref{fig:ablation}): $w_t=0.1$ (block-only) is neutral on Fdataset and Cdataset, where the block geometry already suffices, but costs $0.035$ AUPR on CTDdataset2023 ($p=0.002$), where the tensor view is decisive; $w_t=0.9$ (tensor-only) costs $0.041$--$0.043$ AUPR on Cdataset and Fdataset, and on Ydataset, where the block view is the dominant carry (Figure~\ref{fig:ablation}, \gmcmode{B} AUPR $0.6216$), $w_t=0.9$ is the single worst grid point ($-0.064$, $p=0.002$). Even at that worst point GMC stays competitive with the completion baselines (Fdataset $0.615$ at $w_t=0.9$ vs.\ BNNR $0.600$ and OMC $0.611$; CTDdataset2023 $0.334$ at $w_t=0.1$ vs.\ ITRPCA $0.329$). GMC is insensitive to hyperparameter choice around its operating point and degrades only when an entire view is removed, so the equal-weight center is a robust default rather than a tuned optimum. The flat interior also means GMC transfers to a new dataset without per-dataset tuning.

% Parameter sensitivity is shown in Figure~\ref{fig:param}; key deviations are reported in the narrative below.

\begin{figure*}[tb]
\centering
\includegraphics[width=\textwidth]{figures/fig4_param}
\caption{Parameter sensitivity (AUPR, fresh validation folds): one panel per hyperparameter axis, showing AUPR as a function of the perturbed value (markers) against the center (dashed line) for all four datasets. The equidistant 5-point grid is complete on all four datasets. The interior of every axis is flat; the visible drops are the settings that suppress one of the model's components ($\alpha=0.1$ on the $\alpha$ panel, one-sided $w_t$ on the $w_t$ panel), confirming the reported equal-weight center is a robust default rather than a tuned peak.}
\label{fig:param}
\end{figure*}

\subsection{Robustness to Missing Data}
\label{sec:robust}

The CVa protocol holds out $10\%$ of positive and negative entries. To test how GMC degrades with the degree of missing data (sparsity), Figure~\ref{fig:robust} re-runs the \emph{identical} model under held-out fractions of $5\%$, $20\%$, and $30\%$, keeping the same fold-generation seed (12345) so the $10\%$ point is the reported protocol itself (that point reuses the reported test-fold results); the $5\%$, $20\%$, and $30\%$ points are new masks at the same seed, not the reported folds. As the held-out fraction grows from $5\%$ to $30\%$, AUPR decreases monotonically across all four datasets (Figure~\ref{fig:robust}); the per-dataset values at each mask rate are shown in the figure. Tripling the hidden fraction from $10\%$ to $30\%$ costs $0.06$--$0.09$ AUPR (about $10$--$16\%$ relative) without collapse. Even at $30\%$ masking (roughly one in three associations hidden), GMC retains AUPR of $0.57$--$0.65$, still comparable to several published baselines run at the easier $10\%$ protocol: on CTDdataset2023 the $30\%$-mask GMC ($0.311$) exceeds nine of the eleven baseline AUPRs at $10\%$ masking (only ITRPCA $0.329$ and multiGMF $0.318$ are higher), and on Fdataset and Cdataset it still beats the weaker baselines (four of the eleven on each). The degradation is gradual and roughly linear in the mask rate, not a cliff. GMC is robust to data sparsity, not tuned to the exact $10\%$ protocol.

\begin{figure}[tb]
\centering
\includegraphics[width=\columnwidth]{figures/fig5_robust}
\caption{Robustness to missing data (AUPR vs.\ held-out fraction, 10-fold CVa, fixed model; one line per dataset). All four curves decline monotonically and gently from $5\%$ to $30\%$ masking, with no collapse, showing the model degrades gracefully with sparsity.}
\label{fig:robust}
\end{figure}

\subsection{Case Study: De Novo Prediction and Its Limits}
\label{sec:case}

To probe what GMC can and cannot do beyond the CVa held-out protocol, we consider two extrapolations of the reported configuration.

\paragraph{Strict leave-all-known-out (de novo).} The most demanding test masks \emph{every} known association ($\mathrm{didr}>0 \rightarrow 0$) and asks GMC to recover the association matrix from similarity structure alone. Because the masked matrix is all-zero, the cold-start KNN fill degenerates to zero (its neighbor averages draw on all-zero rows and columns), so the prediction is purely similarity-driven low-rank completion. This is the only configuration in which GMC runs without observed associations, and it isolates the question cleanly, because no other pathway in the model can produce a score independently of the observed entries: the cold-start fill substitutes similarity-weighted averages of its neighbors' labels for missing ones, so it propagates observed associations rather than creating them. We rank all $n_d \times n_s$ entries by predicted score and measure hits@$K$, the number of the top-$K$ that are true (deliberately hidden) associations; a random ranking attains hits@$K \approx K\rho$ at known density $\rho$. The de novo AUPR collapses to the base rate on every dataset ($\approx \rho \approx 0.01$; e.g.\ $0.0086$ vs.\ $\rho = 0.0109$ on CTDdataset2023), and hits@$K$ stays at the $K\rho$ chance level ($0$--$1$ of $10$, $0$--$3$ of $50$; the count fluctuates by $\pm 1$ between runs because the near-flat completion makes the extreme top of the ranking sensitive to thread-dependent float summation in the SVD, but the de novo AUPR is stable). This is a controlled demonstration that GMC's accuracy under CVa is a property of the completion \emph{propagating the observed associations}: the low-rank $+$ similarity prior alone does not rank true pairs above random when no association is observed. The model completes from evidence; it does not infer associations de novo from similarity alone.

\paragraph{High-ranking novel associations (full data).} When the observed structure is available, GMC's extrapolation concentrates on plausible candidates. Running the reported configuration on the full (unmasked) matrix and ranking the pairs with no known association ($\mathrm{didr}=0$), several candidates recur in the top ten of more than one dataset (DB00584--D161900 in three, and DB08799--D607154, DB05271--D168100, DB08802--D607154 among others in two), so the same candidate associations rank stably across datasets built from different similarity views, a weak but non-trivial consistency signal. As a manual follow-up we checked all $40$ top-ten predictions across the four datasets ($29$ of which are unique drug--disease pairs) against the literature via PubMed and CTD: $19$ of the $29$ pairs ($65.5\%$) have direct supporting evidence and a further $6$ ($20.7\%$) are mentioned without established evidence, so $86.2\%$ of the candidates are at least consistent with existing literature, while only $4$ pairs ($13.8\%$) had no evidence found (Table~\ref{tab:case} annotates the displayed top-three rows). These are candidate pairs for follow-up rather than validated findings. Note that this run is the mirror image of the de novo test: the full association matrix is available, so every candidate is the product of completing a nearly complete matrix from abundant observed evidence, not of similarity-based inference. The high hit rate ($65.5$--$86.2\%$) is therefore a statement about ranking quality given observed associations, and it carries no claim about de novo prediction.

% Strict leave-all-known-out (de novo) result: AUPR collapses to base rate, hits@K at chance --- reported in the narrative above.

\begin{table}[tb]
\centering
\caption{Top novel (previously unobserved) drug--disease pairs predicted by the reported GMC configuration on the full association matrix, ranked by predicted score. All pairs have no known association ($\mathrm{didr}=0$). DrugBank and MeSH/CTD identifiers are shown as they appear in the datasets. Ev = manual literature check (PubMed/CTD): $\checkmark$ supporting evidence, $\sim$ mentioned without established evidence, $\times$ no evidence found. The full top-ten lists (forty pairs) are provided in the supporting repository.}
\label{tab:case}
\small
\begin{tabular}{llccc}
\toprule
Dataset & Rank & Drug -- Disease & Score & Ev \\
\midrule
Fdataset & 1 & DB00584 -- D161900 & 0.481 & $\checkmark$ \\
 & 2 & DB01202 -- D267740 & 0.422 & $\sim$ \\
 & 3 & DB00399 -- D166710 & 0.402 & $\checkmark$ \\
Cdataset & 1 & DB08799 -- D607154 & 0.475 & $\checkmark$ \\
 & 2 & DB05271 -- D168100 & 0.472 & $\sim$ \\
 & 3 & DB08802 -- D607154 & 0.469 & $\times$ \\
CTDdataset2023 & 1 & DB00773 -- D603956 & 0.402 & $\checkmark$ \\
 & 2 & DB01169 -- D194070 & 0.400 & $\times$ \\
 & 3 & DB00531 -- D194070 & 0.399 & $\checkmark$ \\
Ydataset & 1 & DB05271 -- D168100 & 0.459 & $\sim$ \\
 & 2 & DB08799 -- D607154 & 0.457 & $\checkmark$ \\
 & 3 & DB01246 -- D607154 & 0.449 & $\sim$ \\
\bottomrule
\end{tabular}
\end{table}

\subsection{Computational Cost}

GMC runs two completion solvers per fold. Let $n=n_d+n_s$ denote the joint block dimension and $\hat r$ the rank cap. The block view is a rank-capped ADMM of $I\le 40$ iterations; each iteration applies one singular-value soft-threshold to the $n\times n$ joint block via randomized SVD, costing $O(n^2\hat r)$ per iteration, hence $O(I\,n^2\hat r)$ overall. The tensor view is a Fourier-domain ALM of $I_t\le 5$ outer iterations over the per-similarity tensors of shapes $(n_d+n_s)\times n_s\times 5$ and $n_d\times(n_d+n_s)\times 2$; each iteration computes a t-SVD (FFT along the slice axis plus a per-frontal-slice SVD), at $O(I_t\,n_3\,n_1 n_2\min(n_1,n_2))$ with full SVD or $O(I_t\,n_3\,n_1 n_2\hat r_t)$ with a rank cap. The cold-start KNN fill is $O((n_d+n_s)\,K\max(n_d,n_s))$ and the rank-normalized fusion $O(n_d n_s)$; both are negligible ($<\!1\%$ of wall time). With the reported full-SVD tensor the tensor view dominates runtime ($\approx\!55\%$ across all four datasets, Table~\ref{tab:runtime}); a rank cap $\hat r_t=200$ cuts the tensor cost $1.7$--$3.4\times$ with AUPR changes within fold-level noise ($|\Delta|\le 0.005$, single fold, and slightly favourable on three of the four datasets) and is the natural lever when speed is preferred over the conservative full-SVD reported here. Both solvers are re-implemented in Python and run end-to-end without external MATLAB dependencies, so the pipeline is self-contained and portable.

Table~\ref{tab:runtime} reports per-fold wall-clock for GMC and the published baselines on the same workstation, both timed on fold 1 with eight pinned threads. GMC runs two completion solvers per fold and pays for it: it is $1.7$--$2.4\times$ slower than BNNR, the lightest completion baseline, on all four datasets, and $2.4$--$2.9\times$ faster than DNMFDDA, the heaviest, while attaining higher AUPR than both (Table~\ref{tab:main_results}). The per-iteration cost is comparable -- GMC's block ADMM is rank-capped with randomized SVD over $40$ iterations, against $300$/$400$ full-SVD iterations in BNNR/DNMFDDA -- so the gap to BNNR is the second solver, not a slower one. With the rank cap $\hat r_t=200$ that gap closes to $11\%$ on the largest dataset (Ydataset) and the lead over DNMFDDA grows to $3.2$--$4.6\times$, at AUPR changes within fold-level noise. The overhead is the price of the two-view design: a single frozen configuration, not tuned per dataset, whose tensor component is decisive on CTDdataset2023 and Ydataset ($+0.045$ and $+0.011$ AUPR) and neutral on Fdataset and Cdataset (Section~\ref{sec:ablation}).

\begin{table}[tb]
\centering
\caption{Per-fold wall-clock (seconds) under CVa on the same workstation. GMC is Python (NumPy/SciPy); baselines are the authors' MATLAB code, run from \texttt{Baseline/}. Every entry is a single fold (fold 1) timed with eight pinned threads (\texttt{OMP\_NUM\_THREADS}=8, \texttt{feature('numThreads')}=8): the block solver is BLAS-thread-sensitive and runs $6$--$10\times$ slower at the $24$ auto-detected threads, so pinning is part of the protocol rather than a best case. The reported GMC configuration uses the full-SVD tensor ($\hat r_t$ uncapped). HGIMC's ChemS+PhS input variant is within $15\%$ of the full-input run shown here; multiGMF's ChemS+PhS variant is $2.6$--$4.5\times$ faster (it fuses one similarity view instead of five drug and two disease views) and is likewise omitted.}
\label{tab:runtime}
\begin{tabular}{lrrrr}
\toprule
Method & Fdataset & Cdataset & CTDdataset2023 & Ydataset \\
\midrule
BNNR      & 8.9  & 14.9 & 25.7  & 73.1 \\
HGIMC     & 2.1  & 3.3  & 2.6   & 10.7 \\
MSBMF     & 6.2  & 8.7  & 13.3  & 44.8 \\
DDA-SKF   & 4.2  & 6.4  & 32.0  & 59.7 \\
OMC       & 6.2  & 9.1  & 16.2  & 42.9 \\
ITRPCA    & 3.9  & 5.3  & 17.3  & 39.8 \\
DNMFDDA   & 46.1 & 62.1 & 178.6 & 356.5 \\
multiGMF  & 2.7  & 3.1  & 10.7  & 20.0 \\
NMF-DR    & 1.5  & 1.3  & 1.3   & 1.2 \\
\textbf{GMC (ours)} & \textbf{18.9} & \textbf{25.2} & \textbf{60.7} & \textbf{131.8} \\
\bottomrule
\end{tabular}
\end{table}

\section{Discussion}

\subsection{Why Completion, Augmented by the Graph, Wins Under CVa}

Under random-entry masking (CVa), each fold leaves a partially observed matrix whose missing entries must be reconstructed, which is precisely the matrix-completion setting where global low-rank methods (BNNR, ITRPCA, DNMFDDA, OMC) lead (Table~\ref{tab:main_results}). Pure label propagation is the wrong tool here: it can only smooth the observed signal and has no mechanism to \emph{invent} structure for randomly masked entries. GMC keeps the winning completion capacity while restoring the local structure that propagation methods exploit. The cold-start-restricted KNN fill supplies a warm prior for novel entities (removing it costs $0.006$--$0.042$ AUPR, Figure~\ref{fig:ablation}, \gmcmode{A}), the block completion embeds that fill in the joint similarity structure, and the tensor view adds a complementary low-rank geometry whose mistakes do not track the block view's. The ablation (Figure~\ref{fig:ablation}) confirms each component earns its keep, and the de novo test (Section~\ref{sec:case}) agrees: with every known association masked, GMC's predictions collapse to the base rate, showing that its accuracy rests on propagating observed associations rather than inferring them from similarity alone. In GMC the distinction is structural: the similarity views are the geometry in which completion happens, not a source of labels, so removing the associations removes the signal itself rather than one input among several.

\subsection{Why Multi-View Fusion Helps}

The block and tensor views see the data differently: the block view uses the \emph{fused} similarities as a regularizer of the joint low-rank projection, while the tensor view keeps each similarity slice separate and projects onto a view-specific low-rank basis. Their errors are decorrelated, so rank-normalized averaging at fixed equal weights exploits this complementarity without overfitting to any single view. The tensor view matters on the datasets where the block geometry falls short (CTDdataset2023 and Ydataset) and is neutral on Fdataset and Cdataset, where the block view alone already reaches the full-model capacity (Figure~\ref{fig:ablation}, \gmcmode{C}); this is the division of labor the fusion exploits. We use fixed equal weights rather than learned ones to avoid overfitting on these small benchmark datasets, and the parameter-sensitivity analysis (Section~\ref{sec:param}) confirms that equal-weight fusion is a robust default across all four datasets.

\subsection{Where the Advantage Is, and Why It Is Not Everywhere}

GMC is significantly ahead of every published baseline on two of the four datasets and statistically indistinguishable from the strongest baseline on the other two. We take the second half of that as informative rather than as a caveat to smooth over.

On Fdataset and Cdataset the tensor view is neutral: ablating it changes AUPR by $|\Delta|\le0.001$, so the reported model reduces in effect to its block view. That is the substrate DNMFDDA also works from, since it receives the fused 5+2 similarities. GMC's mean advantage over it is small ($0.012$ on Fdataset, $0.006$ on Cdataset), smaller than the fold-level standard deviation of either model and with fold-level differences changing sign. Ten folds cannot resolve it, and we do not claim that they have.

On CTDdataset2023 and Ydataset the block geometry underfits. On CTDdataset2023 it is the weakest of GMC's two readouts, with its ablation costing less AUPR than the tensor view's; on Ydataset it is still the dominant component, but the tensor view retains a smaller headroom as well (Figure~\ref{fig:ablation}). The tensor view is decisive there ($+0.045$ AUPR on CTDdataset2023, $+0.011$ on Ydataset, both $p=0.002$), whereas on Fdataset and Cdataset its headroom is essentially zero. A natural account is representational: the block view summarizes the five drug and two disease views with a single mean-fused graph, so structure specific to an individual view is attenuated in it, while the tensor view keeps each view a separate slice and retains it. We offer that as an interpretation and not as a measurement; what we do measure is the consequence, which is that the tensor view has headroom exactly where the block geometry underfits. Whether fusion pays for a given dataset is therefore set by how much the block view alone already recovers on it, not by any choice inside GMC.

Two controls rule out the alternative reading, that the concentration reflects selection rather than data. The reported configuration is frozen before the reported folds are scored (Section~\ref{sec:eval}), and the equal $0.5/0.5$ fusion weights are used unchanged on all four datasets: on Fdataset and Cdataset a tensor-favored setting would have been neutral, while on CTDdataset2023 a block-only setting costs $0.035$ AUPR (Figure~\ref{fig:param}), so the reported numbers cannot be read as having been adapted to each dataset. GMC is not stronger where it was tuned to be; it is stronger where the data leaves room for it to be.

We also do not read the non-significance as a prediction that a better method would fail on Fdataset and Cdataset. The completion family as a whole is compressed into a narrow band there (Table~\ref{tab:main_results}), from BNNR through DNMFDDA to GMC: when the fused block geometry is enough, several different methods land in the same neighborhood. What we cannot do is predict, for an untested dataset, whether the tensor view will pay; the ablation on fresh folds answers that in one run per dataset, and it is the diagnostic we would apply.

\subsection{Comparison with Prior Work}

The multi-similarity family we compare against (MSBMF \citep{yang2020msbmf}, HGIMC \citep{yang2021hgimc}, multiGMF \citep{yang2026multigmf}) was evaluated by its authors under CVa, the protocol we adopt. Under our folds and scoring, the completion members of this family (OMC, ITRPCA, DNMFDDA) are the strongest baselines, and GMC is ahead of all of them. The single-similarity completion baseline BNNR trails the multi-similarity completions on every dataset, confirming the value of the 5+2 views even under completion.

The HN-DRES benchmark \citep{li2024hn_dres} evaluates 28 network-based methods under CVa with 1:1 negative subsampling, where AUPR values are systematically inflated (0.90--0.98). Here we evaluate against the full negative pool, so AUPR values are lower (0.03--0.75); the ranking task is harder and the comparisons more conservative.

More broadly, computational drug repositioning has a long lineage that GMC builds on rather than replaces: heterogeneous-graph inference \citep{wang2013}, label propagation with Gaussian interaction-profile kernels \citep{vanlaarhoven2011}, neighborhood-regularized logistic matrix factorization \citep{liu2016}, neural collaborative-filtering variants \citep{meng2022}, rank minimization with neighborhood-constraint graph integration \citep{liu2025tarm}, and the classic PREDICT formulation of the drug--disease repositioning task itself \citep{gottlieb2011}; recent reviews survey this landscape \citep{huang2025,reda2025}. GMC is methodologically distinct from these families: it treats random-entry masking as a matrix-completion problem and couples global low-rank recovery with local graph structure, rather than relying on propagation, factorization, or kernel smoothing alone. Graph neural networks achieve strong performance under other cross-validation settings, but this work focuses on the matrix-completion paradigm under CVa; integrating GNN modules within the GMC framework is a promising future direction.

\subsection{Limitations and Future Work}

Several limitations merit discussion. First, the two-view design carries a concrete cost--benefit trade-off. The tensor view's full-SVD t-SVD dominates GMC's runtime ($\approx\!55\%$, Table~\ref{tab:runtime}), so GMC costs roughly twice a single completion solver; on Fdataset and Cdataset the tensor view is neutral (no AUPR gain) yet accounts for most of this cost, so its inclusion is pure overhead there. This is the price of a single frozen configuration that is not adapted per dataset: dropping the tensor view on Fdataset/Cdataset would roughly halve runtime without changing AUPR, at the cost of a per-dataset switch we deliberately avoid. On CTDdataset2023 and Ydataset the same cost buys the $+0.045$ and $+0.011$ AUPR that separate GMC from its block-only reduction. A rank-capped tensor ($\hat r_t=200$) cuts the tensor cost $1.7$--$3.4\times$ with AUPR changes within fold-level noise ($|\Delta|\le 0.005$, single fold) and is the natural lever when speed matters; we retain full SVD for the reported numbers as the conservative choice. Second, the cold-start restriction is binary (all-zero vs.\ partially observed); a graded confidence scheme could interpolate. Third, GMC's advantage over the strongest published baseline is demonstrated only on CTDdataset2023 and Ydataset; on Fdataset and Cdataset GMC and DNMFDDA are statistically indistinguishable at ten folds, and a larger fold count or further benchmarks would be needed to separate them. The block-geometry saturation we attribute it to is an interpretation supported by the ablation, not a directly measured quantity. Fourth, GMC is a matrix-completion model and its signal enters only through the observed associations: it cannot rank true drug--disease pairs above chance when no association is observed (Section~\ref{sec:case}), and the case-study candidates are the output of completing a nearly complete matrix rather than of similarity-based inference. Fifth, GMC is evaluated on four public benchmarks; generalization to external, independently curated datasets remains to be tested. Finally, prospective experimental validation of top-ranked predictions would strengthen translational relevance.

\section{Conclusion}

We presented GMC, a drug-repositioning model that treats random-entry masking as matrix completion and combines a global nuclear-norm prior with the local similarity-graph structure that pure completion discards. The same completion is solved in a matrix and a tensor geometry of the identical similarity data and fused on a scale-free axis, with a cold-start-restricted KNN warm start seeding the solver on novel entities. Under 10-fold CVa this single model is significantly ahead of every published baseline on CTDdataset2023 and Ydataset (paired fold-level Wilcoxon, $p=0.002$) and stays ahead of the closest competitor, DNMFDDA, on Fdataset and Cdataset without reaching significance: on those two datasets the block geometry alone already suffices, so fusion has no headroom and the two are statistically indistinguishable at ten folds. The ablation and parameter-sensitivity analyses attribute the gain to the block view's low-rank recovery, the tensor view's complementary geometry, and the cold-start fill, and show the operating point is a robust default rather than a tuned optimum. GMC completes from observed associations, and with every association removed it collapses to the base rate (Section~\ref{sec:case}), so the reported gains are a completion result rather than a demonstration that drug--disease pairs can be inferred from similarity alone. GMC is evaluated on four public benchmarks; transfer to independently curated datasets and prospective experimental validation of the top-ranked predictions remain open.

\section*{Data Availability}

The datasets used in this study are publicly available \citep{yang2020msbmf}. Source code for GMC and all experiments is available at [GitHub URL: to be added upon publication].

\section*{Funding}

This work was supported by [Funding: grant number(s) to be added].

\section*{Conflict of Interest}

None declared.

\section*{Author Contributions}

[Author initials] conceived the method, implemented the algorithms, conducted experiments, analyzed results, and wrote the manuscript.

\section*{Acknowledgments}

We thank the authors of the multi-similarity benchmark methods \citep{yang2020msbmf,yang2019,yang2021hgimc} for making their code and datasets publicly available.

\bibliographystyle{oup-abbrvnat}
\bibliography{references}

\end{document}
